 På \figref{fig:AAUmodel} ses AAU-modellen. 

\begin{figure}[H] 
	\begin{center}
		\includegraphics[width=0.5\textwidth]{figures/cMetode/AAUmodel}
	\end{center}
	\caption{AAU-modellen, der anvendes i projektet. Det ses, hvordan det initierende problem afgrænser det første, brede område, der undersøges i problemanalysen, hvorefter der igen foretages en afgrænsning, når problemformuleringen opstilles.} 
	\label{fig:AAUmodel} 
\end{figure}



\section{Projektets metode}
I dette projekt anvendes HTA core modellen, hvilket afspejles i projektets metode. Målet med HTA core modellen er at analysere og diskutere de essentielle problemstillinger ved en given teknologi, samt at præsentere resultaterne fundet ved analysen og diskussionen \citep{HTAcore}. Ydermere bidrager benyttelsen af HTA-core modellen til en systematisk og bred vurdering af den givne teknologi. Dette giver et evidensbaseret grundlag til beslutningstagning i forhold til eventuel implementation af ny teknologi. \citep{mtvhaandbog2007} \citep{HTAcore}  HTA core modellen anvendes til besvarelse af problemformuleringen, da denne lægger op til en analyse af væsentlige faktorer ved QST set ud fra et sundhedsvidenskabeligt perspektiv. Ved anvendelse af HTA-core inddeles analysen af teknologien i ni forskellige områder \citep{HTAcore}:

\begin{enumerate}
\item Sundhedsmæssigt problem og nuværende brug af teknologien (CUR)
\item Beskrivelse og tekniske karakteristika for teknologien (TEC)
\item Sikkerhed (SAF)
\item Klinisk effektivitet (EFF)
\item Omkostninger og økonomisk evaluering (ECO)
\item Etisk analyse (ETH)
\item Organisatoriske aspekter (ORG)
\item Patient- og sociale aspekter (SOC)
\item Juridiske aspekter (LEG)
\end{enumerate}

Det første analyseområde, CUR, er udført ved benyttelse af en AAU inspireret tilgang. Først indledes projektet med et initierende problem, der lægger op til en bred, alsidig undersøgelse af relevante aspekter for projektet. Herefter følger en problemanalyse, hvis indhold er styret af det initierende problem. Igennem analysen bliver det initierende problem undersøgt o	g afslutningsvis afgrænset, hvormed der kan opstilles en problemformulering. Problemformuleringen definerer det konkrete problem, og forsøges besvaret igennem de resterende analyseområder i HTA core modellen. \\
De første to kapitler (\chapref{introduktion} og \chapref{problemanalyse}), udgør dermed i dette projekt CUR analysen. De resterende områders metode vil blive beskrevet i deres respektive afsnit.  	

Efterfølgende foretages en yderligere inddeling i forskellige emner inden for hvert område. Ud fra emnerne opstilles konkrete problematikker/spørgsmål, der ønskes besvaret. Et område med tilhørende emne og problematik/spørgsmål udgør sammen ét vurderingselement i den samlede teknologivurdering. \citep{HTAcore}. \figref{fig:vurderingselement} viser den beskrevne opbygning af vurderingselementet. 

\begin{figure}[H] 
\begin{center}
\includegraphics[width=0.5\textwidth]{figures/cMetode/vurderingselement}
\end{center}
\caption{Opbygningen af et vurderingselement. Figuren viser, hvordan vurderingselementet samlet er udgjort af de tre elementer; område, emne og problematik.} 
\label{fig:vurderingselement} 
\end{figure}

I dette projekt opstilles vurderingselementerne med udgangspunkt i retningslinjerne i metodehåndbogen for HTA-core modellen. Disse retningslinjer omfatter blandt andet hvilke emner, der bør afdækkes inden for de enkelte områder samt generelle forslag til problematikker. Det fremgår desuden, hvilke andre vurderingselementer, emnet relaterer til. Det fremgår ligeledes, hvilken metode og hvilken type litteratur, der kan være medvirkende til at besvare et givet vurderingselement. \citep{HTAcore} 
Der er ikke altid grundlag for at foretage analyse inden for alle ni vurderingselementer. Ved eksklusion  af et vurderingselement argumenteres der for hvorfor dette element ikke er relevant i den pågældende HTA. \citep{HTAcore} Eksempelvis vil området “Sundhedsmæssigt problem og nuværende brug af teknologien” ikke blive inkluderet i dette projekt, da emnerne og problematikkerne indenfor dette vurderingselement besvares af problemanalysen, som blev lavet forud for problemformuleringen. 
Til besvarelse af den opstillede problemformulering er hver af de tilbageværende otte områder uddybet med emner og problematikker. \textbf{Dette kan eksempelvis ses i afsnit \ref{X}. }
Før besvarelse af hver af disse problematikker argumenteres der for den enkelte problematiks relevans for besvarelsen af problemformuleringen, hvilket sikrer at, kun de problematikker der bidrager til  besvarelse af problemformuleringen besvares i projektet. Efter analyse af hvert område opsummeres vigtige pointer og resultater fra analysen i en delkonklusion.

Afslutningsvis vil besvarelserne på de enkelte vurderingselementer og dermed områderne og problemformuleringen sammenfattes i en syntese indeholdende en diskussion af resultaterne og en konklusion på problemformuleringen. 

\subsection{Litteratursøgning}
Litteratursøgning og -vurdering i projektet tager udgangspunkt i retningslinjerne herfor opstillet i Metodehåndbog for Medicinsk Teknologivurdering udarbejdet af Sundhedsstyrelsen \citep{1}. Da projektet udarbejdes på baggrund af videnskabelig litteratur, er det væsentligt, at litteraturen findes og vurderes ved en systematisk fremgangsmåde, således at problemformuleringen besvares på et veldokumenteret grundlag. Derfor arbejdes der efter en udførlig søgestrategi med en tilhørende søgeprotokol, der sikrer, at udvælgelsen og vurderingen af litteratur forløber overskueligt og ensartet gennem projektet. Ligeledes er formålet med søgeprotokollen at gøre det muligt for interessenter at forstå, hvordan litteratursøgningen er forløbet. \citep{metodehaandbogen}
Til at udarbejde søgeprotokollen i dette projekt, er der opstillet en skabelon, der systematiserer litteratursøgningen. Et uddrag af skabelonen ses i \figref{fig:soegeprotokol}. 

\begin{figure}[H] 
\begin{center}
\includegraphics[width=0.5\textwidth]{figures/cMetode/soegeprotokol}
\end{center}
\caption{Skabelon for projektets søgeprotokol. Det fremgår af skabelonen, at der for hver søgning skal opstilles ét eller flere fokuserede spørgsmål samt inklusions- og eksklusionskriterier. Det skal ligeledes dokumenteres, hvilke databaser, der er anvendt til søgningen. For hver database skal de specifikt anvendte søgeord opstilles, og antallet hits skal efterfølgende fremgå.} 
\label{fig:soegeprotokol} 
\end{figure}

De fokuserede spørgsmål i søgeprotokollen har til formål at bidrage til besvarelsen af problemformuleringen indenfor de ni forskellige områder. Spørgsmålene skal, modsat problemformuleringen, være helt afgrænsede og entydige, således at det er muligt at finde konkret litteratur, og besvare dem præcist. \citep{metodehaandbogen}  
For at afgrænse og sikre relevansen af søgeresultaterne, opstilles inklusions- og eksklusionskriterier for søgningerne. Søgningerne kan eksempelvis afgrænses til kun at indeholde bestemte typer studier, en eller få specifikke sygdomme, en afgrænset aldersgruppe med mere. \citep{metodehaandbogen}\\
Dokumentationen af valgte databaser samt tilhørende specifikke søgeord er en væsentlig del af søgeprotokollen, da databaserne har forskellige sprog til litteratursøgning, hvormed det er vigtigt, at søgningen i en given database er foretaget med anvendelse af søgetermerne specifikke for den valgte database. \citep{metodehaandbogen}
Efter litteratursøgningen vurderes den fundne litteratur med henblik på at udvælge de kilder, der bedst kan besvare de fokuserede spørgsmål og dermed den overordnede problemstilling for projektet. Det er således essentielt at kontrollere, om det, der undersøges i de udvalgte studier, har relevans for de fokuserede spørgsmål eksempelvis ved gennemlæsning af abstract, metode samt studiets resultater. Ligeledes vurderes de enkelte studiers evidensniveau, for at sikre kvaliteten af besvarelserne. \citep{metodehaandbogen} 
Litteraturen i dette projekt inddeles ud fra evidenshierakiet i \cite{metodehaandbogen}. Evidenshierakiet omfatter følgende syv punkter: 

\begin{enumerate}
\item Metaanalyser og systematiske undersøgelser 
\item Randomiserede kontrollerede undersøgelser (RCT’s)
\item Ikke-randomiserede kontrollerede undersøgelser
\item Kohorte undersøgelser
\item Case-kontrol undersøgelser
\item Deskriptive undersøgelser, mindre serier
\item Konsensusrapporter, ikke-systematiske oversigtsartikler, ledere, ekspertudtalelser
\end{enumerate}

Metaanalyser og systematiske undersøgelser er sekundær litteratur og har det højeste evidensniveau. Denne type litteratur er statiske sammenfatninger af primær litteratur med samme afgrænsede problemstilling. \citep{denstoredanske2009} \\
Randomiserede kontrollerede undersøgelser (RCT) er primær litteratur, hvor der foretages en sammenligning af to forsøgsgrupper. Den ene gruppe udsættes for en påvirkning, mens den anden gruppe fungerer som kontrolgruppe. Udvælgelsen af forsøgspersoner foregår tilfældigt. \citep{denstoredanske2009a} \\
Ikke-randomiserede kontrollerede undersøgelser er ligesom RCT’s primær litteratur, hvor to forsøgsgrupper sammenlignes. Ved disse undersøgelser sker udvælgelsen af forsøgspersoner ikke tilfældigt, hvormed evidensniveauet falder da der ikke på samme måde som ved RCT tages højde for bias i forsøgsgrupperne. \citep{denstoredanske2009a} \\
Ved kohorte undersøgelser følges flere forsøgsgrupper over en periode for at undersøge, hvorvidt bestemte eksponeringer har indflydelse på udviklingen af helbredsfænomener, herunder sygdom og død. \citep{metodehaandbogen} \\
I case-kontrol undersøgelser, forsøges det at undersøge forskellige faktorers indflydelse på udvikling af bestemte sygdomme. Dette gøres ved en sammenligning mellem en forsøgsgruppe med den pågældende sygdom og en forsøgsgruppe bestående af raske personer. I modsætning til eksempelvis kohortestudiet følges forsøgsgrupperne ikke over tid, hvormed der ikke kan udføres en opfølgende undersøgelse. Det er således ikke muligt at estimere betydningen af  risikofaktorerne. \citep{denstoredanske2012} \\
Deskriptive studier er studier hvor der foretages analyser til beskrivelse af et fænomen. I modsætning til de andre typer studier påvirkes forsøgsgrupperne ikke. I stedet undersøges nuværende tendenser, eksempelvis med henblik på senere udførelse af et forsøg. \citep{} \\
Fælles for gruppen af konsensusrapporter, ikke-systematiske oversigtsartikler samt ledere og ekspertudtalelser er, at materialet oftest er udtryk for subjektive holdninger, der ikke er underbygget af tilstrækkelige mængder supplerende litteratur, der undersøger området. \citep{metodehaandbogen}  